\documentclass{jletteraddress}

\usepackage{amsmath}
\usepackage[papersize={100mm, 148mm}, margin=0mm]{geometry}

\usepackage{titlesec}

\newcommand{\bra}[1]{\mathinner{\left\langle{#1}\right|}}
\newcommand{\ket}[1]{\mathinner{\left|{#1}\right\rangle}}
\newcommand{\braket}[2]{\mathinner{\left\langle{#1}\middle|#2\right\rangle}}

\sendername{山田　太郎}
\senderaddressa{\hspace*{-10pt}東京都文京区江戸一―二―三}
\senderaddressb{マンションＡ１０１}
\senderpostcode{1000000}

\renewcommand{\baselinestretch}{1.2}

%%% If you want to change the margin of the abstract %%%
%\renewenvironment{abstract}
% {\small
%  \begin{center}
%  \bfseries \abstractname\vspace{-.5em}\vspace{0pt}
%  \end{center}
%  \list{}{%
%    \setlength{\leftmargin}{3mm}% <---------- CHANGE HERE
%    \setlength{\rightmargin}{\leftmargin}%
%  }%
%  \item\relax}
% {\endlist}

\begin{document}

\pagestyle{empty}

\addaddress
  {山田　花子}{様}
  {1000000}
  {東京都文京区本郷二―三―四}
  {}

\newpage

%%%% Change values below to adjust the margin of the body. %%%
\newgeometry{
  top=10mm,
  left=10mm,
  right=10mm,
  bottom=5mm
}

\begin{abstract}
  \footnotesize
  %%% For smaller letters use \scriptsize instead, for example. %%%
  この文章では、新年の挨拶を述べるための手紙である``年賀状''を作る方法について述べる。
  従来の年賀状を作成するソフトウェアは数式やソースコードの埋め込みが貧弱であったが、
  我々はそれを解決す方法として組版ソフト\LaTeX を利用する方法を示す。

  In this paper, we describe how to make a letter to hello new year.
  For conventional softwares that makes those letter,
  it's difficult to write mathmatical expressions or
  programming source code.
  We show how we solve that problem with \LaTeX.
\end{abstract}

%%% If you want to remove the space between the abstract and the body paragraph, adjust by \vspace*. %%%
%\vspace*{-20pt}

\section*{Introduction}
%%% If you prefer a smaller font size for Introduction, add \normalsize before Introduction. %%%

%%% For smaller letters in the body, uncomment \scriptsize here for example. %%%
%\scriptsize

%%% You may adjust the space between the header and the paragraph by %%%
%\vspace*{-8pt}

年賀状は年の始めに互いに送信する手紙のことである。
近年はLINEなどに対抗するため、年賀状に次のような複雑な数式を埋め込みたい
ニーズが存在する。

{\scriptsize
\begin{align*}
  \left|\braket{+}{\varphi_{0 \oplus b, 0 \oplus b}}\right|^2 &= \left\{
\begin{array}{l}
\left|\braket{+}{\alpha\ket{0} + \beta\ket{1}}\right|^2 = \left|\frac{1}{\sqrt{2}}(\alpha + \beta)\right|^2 \\
\left|\braket{+}{\beta\ket{0} + \alpha\ket{1}}\right|^2 = \left|\frac{1}{\sqrt{2}}(\beta + \alpha)\right|^2
\end{array}
\right\} \\
&= \frac{(\alpha + \beta)^2}{2}
\end{align*}
}

\noindent
この文章ではこのような複雑な数式を埋め込むために\LaTeX を
利用した解決方法について述べる。

\begin{flushright}
  \tiny This letter was generated by \LaTeX.
\end{flushright}

%%% If you'd like to add the bibliography, you can write as always. To adjust the space between the body and the references, add %%%
%\vspace*{-10pt}

%%% If you want to lessen spaces between each reference, adjust the space by vspace in the code below. %%%

%\makeatletter
%\renewenvironment{thebibliography}[1]{%
%  \@xp\section\@xp*\@xp{\refname}%
%  \normalfont\footnotesize\labelsep .5em\relax
%  \renewcommand\theenumiv{\arabic{enumiv}}\let\p@enumiv\@empty
%  \vspace*{-10pt} %%% Adjust HERE %%%
%  \scriptsize
%  \list{\@biblabel{\theenumiv}}{\settowidth\labelwidth{\@biblabel{#1}}%
%    \leftmargin\labelwidth \advance\leftmargin\labelsep
%    \usecounter{enumiv}}%
%  \sloppy \clubpenalty\@M \widowpenalty\clubpenalty
%  \sfcode`\.=\@m
%}{%
%  \def\@noitemerr{\@latex@warning{Empty `thebibliography' environment}}%
%  \endlist
%}
%\makeatother

%%% To adjust the font size of the title References, uncomment the following and change the value 9pt as you want. Here BibTeX is used but you may alternatively add references directly as always. %%%
%\begingroup
%\titleformat*{\section}{\fontsize{9pt}{0pt}\bfseries\selectfont}
%\bibliographystyle{unsrt}
%\bibliography{ref}
%\endgroup

\end{document}
